% RESEARCH_DRAFT_SHA256: da28bab463e198290c90a16014d5398a9a86d3d1f4447791f5a4f3749b3f7152
% Layout and source-reference wrappers expanded for the web renderer.
\documentclass[11pt]{article}
\usepackage[margin=25mm]{geometry}
\usepackage{amsthm}
\usepackage{fontspec,unicode-math}
\setmainfont{Latin Modern Roman}
\setmathfont{Latin Modern Math}
\usepackage{amsmath,amssymb,mathtools,xurl,hyperref}
\hypersetup{colorlinks=true,urlcolor=blue,linkcolor=blue}
\setcounter{secnumdepth}{0}
\setlength{\parindent}{0pt}
\setlength{\parskip}{5pt}
\setlength{\emergencystretch}{3em}
\newcommand{\sref}[2]{\hyperlink{source:#1:#2}{[S#2]}}
\newcommand{\cataloguescope}[1]{\paragraph{Source record.}#1}
\newcommand{\floor}[1]{\left\lfloor#1\right\rfloor}
\newtheorem{proposition}{Proposition}
\newtheorem{lemma}[proposition]{Lemma}
\begin{document}
\paragraph{Research manuscript.} OpenAI \texttt{gpt-6-astra}, reasoning \texttt{ultra}; 11 October 2026.
\paragraph{Preserved source body.} The complete catalogue mathematical text and sources follow with layout and source links expanded, including source qualifications and its historical preparation notes. New research begins at the section headed Research attempt.

\section{A causal and decision framework for engaging institutions}\label{statement}

\subsection{Definitions and mathematical statement}

\paragraph{Shared notation.}Standard number systems are $\mathbb N,\mathbb Z,\mathbb Q,\mathbb R,\mathbb C$. Specialized notation and index ranges follow the cited source and the definitions below.


Institutional engagement concerns an organization's participation, allocation of resources, and support for educational implementation. Adoption, instructor training, software maintenance, curricular integration, and evaluation are distinct activities. Open-source software permits source-level reuse and modification under its license; free availability concerns price and does not by itself grant those permissions.
\paragraph{Record scope and qualifications.}
The catalogue title is retained as metadata. Named certificates, contracts, bounds or frameworks in that title are not asserted as results of the original source extract.
This extract contains item 40 followed by extensive summarized and transcribed notes. It is not a single mathematical problem. The entire supplied text is retained literally, including discussion claims that have not been independently verified.

Catalogue screening limitation: Educational assessment, adoption, or study-design discussion; the supplied question does not state a mathematical target with defined assumptions and outcome.

The source snapshot labels this record partially solved; this is catalogue metadata rather than independent verification of a mathematical result or the current completion of a project.
The source text below is rendered literally where necessary; damaged formulas and composite extracts are not silently repaired into a new mathematical claim.
\paragraph{Statement recorded in the supplied source.}
40. Engaging institutions

Summarized Notes

Educational Tools and Open Source vs. Commercial Solutions

The workshop began with a discussion on educational tools like STACK and

WebWork, debating whether to focus exclusively on open-source tools. Participants

distinguished between open-source tools, which can be modified, and freely

available tools, which are not editable. They recognized that while open-source tools

offer control over long-term costs, they still incur expenses related to servers and

expertise. The focus was on ensuring that tools are not only available but also

effectively implemented with proper training and support.

Implementation and Effectiveness

There was a consensus that the success of educational tools depends on their

implementation rather than the tools themselves. Effective use requires thoughtful

integration into the curriculum, considering usability and user community. The

discussion highlighted that diverse assessment methods are needed, and merely

providing tools is not sufficient; critical thinking and training are essential.

Collaborative Problem-Solving and Tool Adaptation

Participants explored the potential of tools designed for collaborative

problem-solving, suggesting that students should be able to pass problems to peers

for continued work. They emphasized the need for technologies that support group

interactions and improve collaborative learning. Additionally, the importance of

adapting tools based on user feedback and ensuring they meet the needs of diverse

learners was highlighted.

Assessment and Data Utilization

The workshop addressed the role of assessments in evaluating student learning,

discussing both intrinsic assessments (where task completion itself verifies correctness) and traditional evaluations. The need for research on the alignment

between digital and traditional assessments was noted, as well as the importance of

integrating meaningful research into teaching practices. Participants stressed the

necessity of understanding how digital tools affect student engagement and learning

outcomes.

Cultural and Contextual Considerations

Participants discussed the cultural aspect of mathematics education, emphasizing

the need to make math relatable and engaging through real-world scenarios. The

conversation also covered the importance of contextualizing online assessments to

address language and cultural differences, and how hybrid methods combining

traditional and technological tools could be beneficial.

Professional Development and Collaboration

The discussion included the need for professional development in using AI and other

technological tools in education. Participants noted the challenges faced by

educators in adopting new methods and stressed the importance of fostering a

collaborative culture in education. Building support networks and addressing

attitudes towards new practices were identified as crucial for effective

implementation.

Future Directions and Research Needs

The workshop concluded with a call for further research into the effectiveness of

educational tools and assessment methods. Participants discussed the need for

longitudinal studies, exploring the impact of technology on different educational

contexts and demographics. They also highlighted the importance of international

collaboration and the need for ongoing evaluation and refinement of educational

strategies.

The workshop aimed to refine ideas into actionable projects and proposals, with a

focus on future collaboration and continued development of educational practices

and tools.

Fully Transcripted Notes

Discussion started with highlighting the two primary tools (STACK and Web work)

but questioned the potential for multiple tools to address challenges in education. The group considered whether to focus on open source tools exclusively and if open

source should be a qualifying criterion.

There was a discussion about the distinction between open source tools, which allow

for code modification, and freely available tools, which may not be editable. The

group aimed to clarify whether "open" means free to use or also includes the ability

to edit and customize the tool.

The total cost of ownership for open source tools was addressed, noting that despite

being free to use, they require servers and expertise, which can be expensive. It was

acknowledged that while open source tools offer control over long-term costs, they

are not completely free, and these costs must be considered by policymakers.

There was a focus on the feasibility and impact of educational interventions,

particularly the use of online tools and assessments. A key point raised by Chris

highlighted the challenge of ensuring these tools are responsive to students' learning

needs and attainment levels. He suggested that the way we use these tools, rather

than the tools themselves, significantly affects their outcomes.

It was emphasized that the tool's effectiveness depends on its implementation and

the policies guiding its use. There was a consensus that simply providing the tools is

insufficient; critical thinking and training on their use are essential.

Two main themes emerged: adoption and implementation. Adoption refers to

whether the tool is used or not, while implementation concerns how the tool is used.

Effective implementation requires considering the theoretical foundations,

user-friendliness, and the community of users.

Participants agreed that technology should be designed to adapt based on feedback

and be supported with appropriate training and resources. The discussion

underscored the need for a comprehensive approach that considers curriculum

views, software usability, and the user community.

It was highlighted that it isn't solely about open source but rather the cost of use

and the implications of maintaining and supporting these tools. The importance of

considering both immediate and long-term costs was emphasized for effective

decision-making in educational contexts.

Furthermore, there was a discussion on the integration of lectureship positions with

curriculum design, emphasizing the importance of creating transferable resources.

The idea is to develop open-source course packs, such as those for linear algebra,

that instructors can download and use. These resources should not only be

well-packaged for use but also designed for sharing and community collaboration,

allowing for feedback and continuous improvement. A key point raised was the need for a two-way conversation in resource sharing.

Instead of a top-down approach, where a package is distributed for everyone to use,

the focus should be on building a collaborative loop. This involves educators

contributing to and refining shared resources.

Additionally, the workshop highlighted the challenge of authoring questions, which is

time-consuming for educators. There was a discussion about the interoperability of

educational content across different learning systems. This includes the potential for

importing and adapting course materials from one platform to another, ensuring that

content is reusable and efficient.

There was a call for both technological and community-based solutions to facilitate

the sharing of educational resources. The goal is to improve the quality and volume

of shared content, thereby saving time and enhancing the overall educational

experience.

Another discussion on accessing education data emerged and emphasized the

importance of tailoring educational tools and data collection to different contexts to

motivate learners effectively. This contextual approach ensures that data is

representative of diverse environments, aiding in comprehensive analysis.

Participants highlighted the need for large datasets to train effective models.

Collaboration with institutions is essential to gather socio-demographic information,

which can enhance the utility of data for various purposes. A key question raised was

the broader objectives of collecting combined data sets and the types of questions

such data could help answer.

One significant barrier to technology adoption in education is the lack of adaptability

to individual learner levels. Technologies often fail to identify specific areas where

students struggle, unlike human teachers who can provide personalized guidance.

Addressing this gap could involve using data to adapt educational technologies to

meet individual learning needs more effectively.

Overall, the discussion underscored the necessity of actionable data to identify and

address learning gaps, enhancing the adaptability of educational tools to support

student success.

The concept of intrinsic assessment was discussed, particularly in the context of

problem-solving and modeling activities. Intrinsic assessment occurs when the

completion of a task inherently demonstrates its correctness. For instance, in coding,

the functionality of the code serves as its own assessment-if it works, it meets the

required standards. This self-assessing nature is valuable but should be one of many

tools in an educational portfolio. Discussion on authenticity and functionality highlighted that for an artifact to be

considered authentic, it must function correctly. This idea challenges the traditional

"us versus them" model of assessment, where an external party evaluates the work.

Instead, the artifact's ability to perform its intended function serves as a measure of

its authenticity and correctness.

Broadening assessment tools discussions underscored the importance of having a

diverse set of assessment tools. Intrinsic assessments, while valuable, cannot stand

alone. Educators should incorporate various methods to ensure comprehensive

evaluation and support student learning.

There was also a discussion on the need to research the effectiveness of new

educational tools, such as STACK, in enhancing student learning. Concerns were

raised about potential unintended consequences of these innovations. It was

suggested that thorough testing and research are necessary to understand their

impact fully and to address any negative outcomes.

The interaction between students and educational tools was another key topic. The

importance of structured time and focused engagement was emphasized to prevent

students from rushing through tasks without understanding. The debate about the

quality of online math practice compared to traditional methods was also addressed,

with suggestions that research could help validate the effectiveness of online tools.

There was a consensus on the need for diverse assessment methods, careful

implementation of educational innovations, and thorough research to ensure these

tools positively impact student learning. The discussions highlighted the complexities

of modern education and the necessity of a multifaceted approach to teaching and

assessment.

Incorporating math education research at the development stage of technologies can

provide valuable feedback to improve teaching and learning. One key area needing

research is the development of teachers' content knowledge. For instance,

understanding how to effectively teach fractions and identifying common student

mistakes can be challenging. Technology can help by collecting and analyzing data

on student performance, which can then be used to inform teacher training and

improve instructional methods.

In the Kenyan context, the shift from a summative to a formative assessment

approach under the Competency-Based Curriculum (CBC) highlights the need for

better utilization of assessment data. By analyzing data from formative assessments,

the government can provide feedback to teachers, helping them address specific

areas of student weakness. This approach can enhance both individual and national

education outcomes. Research should also focus on the specific features of assessment tools that support

student engagement with mathematical ideas. Understanding how feedback is

structured and presented can be crucial. Qualitative research, such as interviewing

students about their experiences, can provide insights into what supports or hinders

their learning. This information can guide the design of more effective feedback

mechanisms, ultimately improving student learning outcomes.

Online assessments often fail to connect with students due to contextual differences.

One significant issue is the language used in these assessments, which is

predominantly English. The expectations for how responses should be input can be a

barrier, especially if the student's way of expressing themselves isn't aligned with

conventional standards. Educators who know their students well can often infer their

intended meaning, but this nuance is lost in automated online assessments.

There is a need to explore ways to bridge this gap and make online assessments

more context-sensitive. One suggestion is to incorporate tools that allow for some

level of interpretation or personalization based on the student's context. Additionally,

innovative approaches to teaching and assessment that consider the specific context

and needs of students should be developed. Hybrid methods combining traditional

and technological tools could be beneficial.

Mathematics should not just be viewed as a subject but as a cultural element that

influences various professions. The discussion highlighted that individuals exposed

to mathematical thinking from an early age tend to excel in their fields, even if those

fields are not directly related to mathematics. This cultural aspect of mathematics

helps individuals develop better problem-solving skills and analytical thinking, which

are valuable in any profession.

The example of using a golf ball to teach mathematics illustrates the importance of

making math relatable and applicable to real-world scenarios. This approach can

change students' perceptions of mathematics and make it more engaging and

relevant to their lives and future careers.

Discussions here emphasized the importance of contextualizing online assessments,

using hybrid methods, and promoting the cultural aspects of mathematics to enhance

learning outcomes. These strategies can help bridge the gap between students'

understanding and conventional assessment methods, ultimately fostering a deeper

appreciation and proficiency in mathematics.

Participants discussed the potential of designing educational tools that facilitate

collaborative problem-solving. One idea presented was a tool allowing students to

start solving a problem individually and, if they get stuck, pass it on to a peer who

continues the work. This approach would enable students to learn from each other's problem-solving methods. The discussion emphasized shifting from technology

designed for individual use to technology that supports group interactions. This shift

could enhance collaboration, particularly in the context of competency-based

curricula and 21st-century skills.

Collaboration in mathematics goes beyond group work; it involves students sharing

and building on each other's ideas. Technologies that support this kind of interaction

can foster deeper collaboration and improve learning outcomes. Participants

explored the idea of involving students in content creation, not just as consumers.

This approach could address language barriers and content accessibility, making

learning materials more relevant and authentic. The discussions addressed the

challenge of ensuring that student feedback and answers in assessments are

authentic. Participants discussed the need for reliable electronic tools that accurately

reflect students' understanding and performance.

The discussion highlights a key issue: the gap between current technology use and

the experience of educators who may not have been trained in modern tech-based

teaching methods. The focus needs to be on a holistic approach that considers the

entire educational system, including policymakers, educators, and students. There is

an observed resistance or lack of familiarity with new methods among educators, not

necessarily due to opposition but because they have not been exposed to or trained

in these modern approaches. This suggests the need for a shift in training programs

for future educators to better integrate contemporary practices. The conversation

also emphasized the importance of addressing attitudes and values in educational

change. Successful implementation of new practices, whether technology-based or

not, requires attention to the attitudes of those involved. This means incorporating

these aspects into the design and deployment of educational initiatives to ensure

effective adoption and application.

There were discussions revolving around how to effectively build and sustain a

community around educational technologies like STACK, ensuring high adoption and

ongoing development. A major concern is how educators using STACK can interpret

the data analytics it provides, especially since not all users have a background in

statistics. The goal is to simplify this data so that educators, regardless of their

statistical expertise, can easily understand and apply the insights to address specific

issues their students may face.

The question posed is how to make the analytics from tools like STACK more

accessible and useful for educators. It is essential to explore ways to automate or

simplify the process of interpreting and sharing insights from these tools. Additionally,

understanding how these tools impact different types of engagement-emotional,

cognitive, and behavioral-is important. This includes examining whether these tools

affect engagement levels differently and using this understanding to guide future

improvements. In summary, the discussion sought to address how to enhance the usability of

educational tools and the analytics they provide, focusing on improving their

accessibility for educators and understanding their impact on student engagement.

Again, the discussions highlighted several key issues around communication and

student engagement in educational settings. One notable point was the impact of

transitioning to digital tools on student-instructor relationships. An example

mentioned was about how a professor shared that switching to online homework

submissions reduced their familiarity with student names, demonstrating how

technology can affect personal interactions.

Discussions emphasized the importance of maintaining student interaction, even

when integrating new technological tools. It was argued that while digital tools can

enhance learning, they should not replace face-to-face engagement. This balance is

critical to ensuring students feel heard and supported.

Participants suggested that instead of relying on a single comprehensive tool, a suite

of complementary tools might better address various educational needs. This

approach allows instructors to choose the most appropriate tools for quizzes, group

work, assessments, and content delivery based on their specific class context.

One proposed strategy was using tools to foster student interaction and

collaboration, where students receive additional points for helping their peers. This

method encourages active participation and peer support, contributing to a more

dynamic learning environment.

The discussion concluded with a call for a stable, long-term platform for educators to

share and receive feedback on the effective use of technological tools. This platform

could help educators adapt and improve their teaching strategies, ensuring that

technology enhances rather than detracts from the learning experience.

Discussions further highlighted the effectiveness of structured pedagogical activities

for teachers. By following well-designed activities step-by-step, even less

experienced teachers can see improvements in teaching and learning outcomes.

However, this approach may limit opportunities for innovation and creativity, which

could be a drawback for confident teachers looking to enhance their sessions further.

A key topic was the concept of "collaboratively invented mathematics," where

students use digital tools to collaboratively discover mathematical concepts, such as

the Taylor series. This method shifts from traditional pedagogy to a more engaging,

exploratory approach, allowing students to invent mathematics that historically took

centuries to develop. There was a strong emphasis on embedding meaningful research into teaching

practices. This includes designing assessments with digital tools to evaluate

students' understanding effectively. The integration of research can inform

educational innovations and improve teaching methodologies.

Participants also discussed the importance of international collaboration in

mathematics education. Countries interested in adopting these innovative teaching

methods need support to integrate and implement them effectively. The potential for

using online tools to facilitate these collaborations was considered crucial for

broadening the impact of these educational innovations.

It was then concluded from this discourse that structured pedagogy can significantly

improve teaching outcomes, but there is a need to balance this with opportunities for

teacher innovation. The collaborative invention of mathematics and embedding

research into teaching practices were highlighted as promising approaches. Global

collaboration and effective implementation of these methods are essential for their

success.

A participant had mentioned the need to discuss the positive uses of AI in teaching

and its potential benefits. Participants highlighted the importance of professional

development for lecturers and teachers to effectively use AI tools, including both

pre-service and in-service training.

A concern was raised about why only a few lecturers consistently use new

technological tools while others do not. The discussion also explored the support

available for African institutions wishing to adopt technology in teaching. Building

support networks and fostering collaborative learning were identified as crucial

elements.

The conversation noted that education systems often promote individualism over

collaboration. This mentality persists into higher education and research, making

collaboration challenging. Encouraging a culture of helping and sharing knowledge

from a young age was seen as vital.

Also, it was mentioned that some teachers appreciated innovative teaching methods

that do not rely on technology but expressed concerns about time constraints. They

feared that creative teaching methods might reduce the amount of content covered

during class. The group questioned ways to balance innovative teaching with

curriculum requirements, aiming to inspire students to explore concepts

independently.

One issue discussed was the alignment between final exam results and outcomes

from online assessments. The concern is whether traditional exams provide the

same results as digital formative assessments. This disparity could be an interesting research question, particularly in contexts like the Open University, where online

assessments are prevalent. Understanding how long-term use of digital tools affects

mathematical communication and writing could be another research avenue.

A key point raised was the relationship between formative assessments and

traditional examinations. There is interest in researching how these different forms of

assessment align with each other, especially if one is digitized and the other is not.

This could reveal important insights into the effectiveness and consistency of various

assessment methods.

Another discussion topic was the role of digital tools in enhancing or hindering

mathematical communication. The group considered how these tools impact

students' abilities to communicate mathematical ideas effectively. There was a

suggestion to explore ways to leverage student collaboration to improve

communication skills, potentially by rewarding students for explaining concepts to

peers.

The conversation also touched on how improving collaborative learning can

simultaneously enhance mathematical communication skills. Encouraging students

to work together and explain their reasoning can create a virtuous cycle of improved

communication and understanding. This approach could be beneficial in fostering

both collaboration and competency in mathematics.

Finally, participants highlighted the need to understand different levels of

mathematical education. From high school students aiming for basic competency to

those pursuing careers in mathematics, it's important to consider how various tools

and methods support different educational goals. This broader understanding can

help tailor educational strategies to meet diverse student needs.

The math education researchers discussed the importance of building capacity

among mathematicians who currently teach but may lack certain skills. This involves

engaging more individuals in math education research beyond just the math

education researchers.

Participants considered conducting a research project aimed at improving the quality

of math tasks. This includes identifying the best and worst tasks and determining

where new tasks should be developed. Research on individual tasks can help

pinpoint those that are most effective. There are challenges in designing tasks for

certain areas of mathematics, such as abstract algebra and geometry. The

discussion covered the need to create effective tasks that go beyond simple

calculations and consider the specific content of each course. An idea was proposed

to identify a set of principles for designing good tasks that promote learning,

regardless of the course context. For instance, out of 200 derivative problems,

finding the top 20 that best promote learning. The workshop emphasized the importance of maintaining the human element in

assessments. Positive reinforcement can generate new ideas and support students

effectively. One participant shared their teaching approach, which involves assigning

seminar topics to groups of students. These groups work on their topics throughout

the semester and present their findings, fostering collaboration and deeper

understanding.

During the workshop, discussions focused on the integration of education technology

and its impact on student learning. One key point raised was the need to assess

whether these technologies are genuinely beneficial or potentially harmful. Although

initial plans to collaborate with a math education researcher were not realized in

time, the intention is to pilot the technology with first-year students and conduct a

follow-up assessment in the second year.

Participants emphasized the importance of understanding the distinct roles education

technology plays in different institutional contexts, such as CalTech versus other

universities. This understanding is crucial for determining the effectiveness of such

technologies in enhancing learning experiences.

Another significant idea was the implementation of longitudinal studies to track

student progress over several years. These studies could help identify best practices

and measure the long-term impact of education technologies on learning outcomes.

For example, tracking the same cohort of students through a four-year degree

program could reveal valuable insights into their learning journeys.

The workshop also highlighted the potential to investigate specific issues, such as

gender disparities in STEM fields. By comparing data from different universities and

contexts, researchers could analyze how online tools and other interventions

influence retention rates and learning experiences for different student

demographics.

In conclusion, the discussions underscored the need for rigorous educational

research to identify effective practices and understand how various factors influence

student learning across different contexts.

Discussions highlighted the tendency to treat students as a homogenous group,

which overlooks individual progress. Mary raised a point about the uniform speed of

student progress enforced by current assessment systems. Unlike learning to drive

in the UK, where individuals take their driving test when ready, school exams are

scheduled uniformly for all students. This system's rigidity does not account for

individual readiness and progress. The conversation explored the potential of

electronic assessment tools to transform not only learning but also assessment

systems. The current system, rooted in historical practices, necessitates uniform

exam schedules. Electronic tools, however, offer the flexibility to tailor assessments to individual progress, potentially leading to a more personalized and effective

education system.

Mary's quote about play sparked a discussion on the nature of compulsory

participation. True play requires freedom-emotional, economic, and choice

freedom. Compulsory education systems often lack these freedoms, making

participation feel forced. The group pondered whether new tools could introduce

more freedom and playfulness into learning, thereby enhancing engagement and

effectiveness.

The discussion extended to the broader implications of changing educational

systems. It was suggested that learning could be more context-specific and playful,

integrating disciplines in meaningful ways. For instance, learning math through

historical contexts could make it more relevant and engaging for students.

A key concern was whether stakeholders-students, educators, institutions, and

societies-are ready for such a transformation. The readiness in terms of attitude,

capacity, and resources was questioned, especially considering the challenges faced

by educational systems in regions like Africa. The discussion concluded with a call to

assess the readiness and willingness of all stakeholders to transition from traditional

to technology-integrated assessments.

The workshop highlighted a critical distinction between commercial and open-source

solutions, particularly concerning future cost implications. Participants debated

whether the focus should be on Open Access or open-source terms, considering

their impact on accessibility and contribution rights.

A key point raised was the risk of relying on commercial software that might become

costly or inaccessible if terms change. In contrast, open-source software offers more

stability and control, allowing modifications and reducing dependency on external

vendors.

The discussion also touched on the need to understand how open-source principles

could inform educational software choices. The idea is to draw from the open

community's experiences to determine what makes software truly open and

sustainable.

Participants expressed interest in identifying qualifying criteria for evaluating different

software options. The upcoming sessions will aim to clarify these criteria, explore

alternative platforms, and discuss their functionalities to make informed decisions.

Overall, there is a call to explore how open-source and open-access principles can

better serve educational institutions and to determine the most suitable approach

based on their needs and trade-offs. Another discussion highlighted the tension between competitive academic systems

and the need for collaboration. There is a concern about how to foster collaboration

from a young age within a system that traditionally emphasizes competition. This

raises questions about how competitive academic structures can adapt to support

collaborative learning.

It was noted that shifting teaching approaches might not require technology but a

change from a strict curriculum focus. Teachers accustomed to curriculum-based

instruction may struggle to integrate new methods. The question arises whether

these new approaches can fit within the existing curriculum or if they require a

complete overhaul.

The conversation addressed the need to adapt digital tools for students with special

needs. Ensuring that digital educational resources are accessible to all learners is

crucial, and there is interest in how these tools can be modified to meet the needs of

individuals with disabilities.

A question was raised about whether students could receive certification for

completing modules from open educational resources outside traditional institutions.

This discussion explores the potential for recognizing and certifying informal or

self-directed learning experiences.

Recent evidence suggests that structured pedagogy, which includes detailed lesson

plans and specific instructional strategies, improves teaching and learning at

foundational levels. The inquiry is whether similar structured approaches could

enhance education at higher levels, combining structured methods with more

advanced pedagogical strategies.

The discussion highlighted the importance of exploring alternative assessment

methods beyond digital tools. Emphasis was placed on incorporating human

interactions and experiences, which are often difficult to quantify. These

assessments should inspire and motivate rather than merely evaluate.

There is a need for online tools that not only assess but also engage and inspire

users. This involves considering how institutions can be involved in data sharing

agreements and fostering better engagement at an institutional level, rather than

focusing solely on individuals.

As the session concluded, there was a brief discussion on record-keeping and the

need for capturing high-speed data. The final points stressed were the importance of

integrating motivational aspects into assessments and the need for institutional

involvement in data sharing.

The workshop aims to explore and develop a variety of significant and

well-considered ideas. Participants will engage in open-ended discussions and collaborative activities to refine these ideas. The goal is to convert these ideas into

actionable projects, such as grant proposals or collaborative efforts.

The moderators and organizers will review the collected ideas and organize them

into key topics for group discussions scheduled for tomorrow. This process will

involve multiple cycles of group work to generate viable projects.

By the end of the workshop, participants are expected to develop detailed plans and

potential projects. However, given the scope of the topics, the workshop will primarily

serve as a starting point, with continued work beyond the event.

Participants should use the current session to propose and refine ideas they are

interested in. The workshop will provide a foundation for future collaboration, with the

understanding that comprehensive solutions will evolve over time.

A report summarizing the workshop outcomes will be available, detailing the

proposed topics and next steps. Participants are encouraged to bring forward any

new ideas or questions they have. \href{http://aimath.org/pastworkshops/oerassessmentproblems.pdf}{[S2]}
\paragraph{Source status.}
The supplied catalogue labels this source record partially solved; that metadata was checked against the supplied snapshot on 2026-10-11. The label is preserved for provenance and is not an independent verification of present open status or mathematical correctness.
\subsection{Short English statement}
Preserve the institutional-engagement prompt and its bundled workshop discussion about educational tools, implementation, and research.
\subsection{Sources}
\begin{itemize}
\item[S1] Ulam AI and the UnsolvedMath Contributors, UnsolvedMath: A Curated Collection of Open Mathematics Problems, catalogue entry. Catalogue curation and metadata: CC BY 4.0. The source extract is retained; definitions, English summaries and qualifications are editorial additions.
\url{https://huggingface.co/datasets/ulamai/UnsolvedMath}.
\item[S2] AIM Workshop Problem Lists, source linked by the catalogue.
\url{http://aimath.org/pastworkshops/oerassessmentproblems.pdf}.
\end{itemize}
\paragraph{Preparation and review.}
The original source and UnsolvedMath contributors retain attribution for the source extract. Codex (GPT-6; exact serving revision not exposed) assisted with definitions, English summaries, scope notes and formatting. No new human mathematical authorship or independently verified solution is asserted. Catalogue identity checked against the complete supplied registry and current repository. Distinct upstream records, including repeated questions and broad or damaged extracts, are retained for complete coverage of the source CS category. A repeated record is not a new independent conjecture. Complete mathematical-equivalence and current-literature review has not been performed. Preparation and final review follow contribution prompts C and D, with source gaps disclosed. Structural and rendering checks do not establish a new theorem or open status.
\label{end:statement}

\section{Research attempt}
\subsection{Literature, source context and scope}
The primary AIM notes describe a brainstorming session. The original excerpt and all catalogue qualifications remain above. They do not specify a mathematical theorem or a complete empirical study. The formal model below is an explicitly limited research route introduced for this attempt, not a reconstruction of a missing original conjecture. The supplied catalogue's earlier research summary was inspected as context; its named framework is not adopted as an established solution. The elementary result below is not claimed novel, and no educational effectiveness data were collected.

\subsection{Attempt: complementary institutional commitments}
This source record contains item 40 followed by a long composite workshop transcript, all preserved above. It has no single formal success condition. I tested only one institutional-coordination mechanism: in a finite directed network, an institution becomes active once at least its declared threshold number of incoming neighbors are active. Initial seeds are active unconditionally, and activation is irreversible. This is a hypothetical threshold model, not a conclusion about actual institutions or the rest of the transcript.
\begin{proposition}
Repeated threshold activation reaches the unique least closed active superset of the seeds after at most $|V|$ strict growth steps. The resulting active-set size as a function of seed set need not be submodular, even on three vertices.
\end{proposition}
\begin{proof}
For seed set $S$, set $S_0=S$ and
\[
S_{t+1}=S_t\cup\{v:|N^-(v)\cap S_t|\ge\theta_v\}.
\]
This sequence only grows. Each strict growth adds at least one of the finitely many vertices, so after at most $|V|$ such steps it reaches a fixed point. If $C$ is any closed set containing $S$, induction gives $S_t\subseteq C$: any vertex meeting its threshold in $S_t$ also meets it in $C$, hence belongs to $C$. The fixed point is therefore contained in every such $C$, proving it is the unique least one.

Take edges $u\to w$ and $v\to w$, thresholds $\theta_w=2$ and $\theta_u=\theta_v=1$. The last two vertices have no incoming neighbors, so they can activate only as seeds. Let $f(S)$ be the final active-set size. Then $f(\varnothing)=0$, $f(\{u\})=f(\{v\})=1$, and $f(\{u,v\})=3$. Submodularity would require $f(\{u\})+f(\{v\})\ge f(\{u,v\})+f(\varnothing)$, but this says $2\ge3$. Thus complementary commitments violate diminishing marginal gains.
\end{proof}
\subsection{Verification and remaining gaps}
Direct closure iteration on all eight seed subsets checks the four values used in the obstruction. The failed route is applying an independent-increment or submodular-benefit selection guarantee merely because adding institutions cannot reduce participation. Monotonicity does not imply diminishing returns. Actual commitment rules, incentives, costs, evidence-sharing arrangements and educational outcomes are not supplied, and none is inferred from this graph example. The full composite institutional and educational agenda remains unresolved; no original-target completion is claimed.

\section{References and source attribution}
American Institute of Mathematics, workshop discussion notes on open educational resources and assessment, \url{https://aimath.org/pastworkshops/oerassessmentproblems.pdf}. The 16-page primary PDF was opened on 11 October 2026 and its numbered brainstorming item was checked against the supplied record. The source supports the provenance and agenda, not the additional modeling assumptions or an empirical educational conclusion. Catalogue and contributor attribution are retained in the reproduced Sources section. The worked mathematics is proved here without reliance on an unverified external theorem.


\paragraph{Integration review.} Codex (GPT-6; exact serving revision not exposed) reviewed the written arguments and expanded layout and source-reference wrappers. Repository metadata and the production web renderer were checked using Python 3.13 and JavaScript with MathJax 3.2.2. This model-assisted review is not independent human proof verification; no PDF compilation or formal proof checking is asserted.

\section{Completion Estimate}
\noindent\textbf{COMPLETION: 0\%}
No substantiated new progress toward the entire original target is claimed. The worked results and failed route are delimited above; a known result, illustrative model, or auxiliary obstruction is not counted as original-target completion.

\section{Final status and attribution}
\textbf{Attempt status: unresolved.} This describes this research attempt, not a new ruling on the source catalogue's status or on external literature claims.
The mathematics and this self-audit were generated by OpenAI \texttt{gpt-6-astra}, reasoning \texttt{ultra}, on 11 October 2026. Codex (GPT-6; exact serving revision not exposed) orchestrated the batch and contributed formatting/repository review. No human mathematical authorship, independent human proof checking, or formal verification is claimed. The historical preparation paragraphs in the reproduced catalogue body concern that body only.

Contribution prompt A was applied to the complete definition and supplied source record, followed by the adversarial checks of prompt B. Local reproduction files are \texttt{batch008/checks.py}, \texttt{computations.json}, and the per-ID audit JSON. They were not public links at preparation time. Numerical checks are finite and do not prove unrestricted results. No LaTeX engine was found locally; compilation was not performed by this research agent.

\end{document}
